\documentclass[11pt]{beamer}
%\usetheme{Warsaw}
\usepackage[utf8]{inputenc}
\usepackage[portuguese]{babel}
\usepackage[T1]{fontenc}
\usepackage{amsmath}
\usepackage{amsfonts}
\usepackage{amssymb}
\usepackage{graphicx}
\author{eu}
\title{teste}
%--------------------------------------------------
% Lista de temas(testei um por um haha)

%\usetheme{default}
%\usetheme{AnnArbor}
%\usetheme{Antibes}
%\usetheme{Bergen}
%\usetheme{Berkeley}
\usetheme{Berlin}
%\usetheme{Boadilla}
%\usetheme{CambridgeUS}
%\usetheme{Copenhagen}
%\usetheme{Darmstadt}
%\usetheme{default}
%\usetheme{Dresden}
%\usetheme{Frankfurt}
%\usetheme{Goettingen}
%\usetheme{Hannover}
%\usetheme{Ilmenau}
%\usetheme{JuanLesPins}
%\usetheme{Luebeck}
%\usetheme{Madrid}
%\usetheme{Malmoe}
%\usetheme{Marburg}
%\usetheme{Montpellier}
%\usetheme{PaloAlto}
%\usetheme{Pittsburgh}
%\usetheme{Rochester}
%\usetheme{Singapore}
%\usetheme{Szeged}
%\usetheme{Warsaw}

%--------------------------------------------

% Mudar a cor do tema.(decidi não mexer)

%\usecolortheme{albatross}
%\usecolortheme{beaver}
%\usecolortheme{beetle}
%\usecolortheme{crane}
%\usecolortheme{dolphin}
%\usecolortheme{dove}
%\usecolortheme{fly}
%\usecolortheme{lily}
%\usecolortheme{orchid}
%\usecolortheme{rose}
%\usecolortheme{seagull}
%\usecolortheme{seahorse}
%\usecolortheme{whale}
%\usecolortheme{wolverine}

%-----------------------------------------------------------------

\setbeamertemplate{footline} % Remover o rodapé de todas as páginas
%\setbeamertemplate{footline}[page number] % Trocar o rodapé pelo nr de páginas
%\setbeamertemplate{navigation symbols}{} % Remove todos os símbolos de progresso dos slides


\usepackage{graphicx} % Allows including images
\usepackage{booktabs} % Allows the use of \toprule, \midrule and \bottomrule in tables

%----------------------------------------------------------------------------------------
%	Capa
%----------------------------------------------------------------------------------------

\title[Estatística Aplicada a Geopolítica]{Análise do Impacto Socioeconômico de Serviços Inovadores que Competem com Monopólios ou Cartéis} % o que está dentro dos [] aparecerá em todos os slides e dentro das {} aparecerá na capa
\author{Jayme Gomes dos Santos Junior} 
\institute[Universidade Federal do Paraná] 
{
  \inst{}%
  Paulo Afonso Bracarense Costa\\
  Orientador
  }
\date{}
%\textit{jaymegomesdossantosjunior@gmail.com} % e-mail(decidi não deixar)
%}

\logo{\includegraphics[width=3cm]{logo-pet-azul}}
%-------------------------------------------------------------------------------------

\usepackage{Sweave}
\begin{document}
\Sconcordance{concordance:siepe.tex:siepe.Rnw:%
1 85 1 1 0 132 1}


\begin{frame}
\titlepage % Mostra a capa como primeiro slide
\end{frame}

%\begin{frame}
%\frametitle{índice} % Sumário
%\tableofcontents % Se forem acrescentadas sessões e subsessões, com esse comando, elas serão colocadas automaticamente no sumário
%\end{frame}

%----------------------------------------------------------------------------------------
%	Slides
%----------------------------------------------------------------------------------------

%------------------------------------------------
%\section{Motivação} 
%------------------------------------------------

%\subsection{Teoria de Rifkin} 

%------------------------------------------------

%\begin{frame}
%\frametitle{Motivação - Terceira Revolução Industrial}
%\pause
%\begin{itemize}
%\item Teoria de Rifkin
%\pause
%\item Eclipse do capitalismo como conhecemos
%\pause
%\item Mudança de paradigma econômico
%\pause
%\item Terceira revolução industrial
%\end{itemize}
%\end{frame}

%------------------------------------------------

\begin{frame}
\frametitle{Objetivo}
%\pause
%\begin{itemize}
%\item Teoria de Rifkin
%\pause
%\item Eclipse do capitalismo como conhecemos
%\pause
%\item Mudança de paradigma econômico
%\pause
%\item Terceira revolução industrial
%\end{itemize}
\begin{block}

A pesquisa tem como objetivo encontrar evidências
de que o fenômeno da mudança de paradgma econômico
apontado por Rifkin está acontecendo no Brasil.
\end{block}
\end{frame}

%------------------------------------------------
%\section{Materiais de Apoio}
%------------------------------------------------

\begin{frame}
\frametitle{Metodologia}
\begin{itemize}
\item Monitoramento de tweets no período de 17/04 a 31/08 de 2017;
\pause
\item Coleta dos tweets e extração de seus conteúdos;
\pause
\item Tratamento e análise dos dados.
\end{itemize}
\end{frame}

%------------------------------------------------

\begin{frame}
\frametitle{Análise de Sentimento}
\begin{itemize}
\item Gráfico aqui;
\pause
\item Mais um gráfico aqui;
\end{itemize}
\end{frame}

%------------------------------------------------
%\section{Ferramentas}
%------------------------------------------------

%\subsection{Mendeley}

\begin{frame}
\frametitle{Notas de Conclusão}
\begin{itemize}
\item Desafios;
\pause
\item Evidência;
\pause
\item Futuro.
\end{itemize}
\end{frame}

%------------------------------------------------

%\subsection{R}

%\begin{frame}
%\frametitle{Mineração de Texto}
%\pause
%\begin{itemize}
%\item Tweeter - API e pacote R
%\pause
%\item Pacote Selenium para o Google Play
%\pause
%\item Ferramentas de manipulação de texto
%\pause
%\item Análise posterior
%\end{itemize}
%\end{frame}

%------------------------------------------------
%\section{Desafios}
%------------------------------------------------

%\begin{frame}
%\frametitle{Hiato}
%\begin{itemize}
%\item Desorganização
%\item Tempos difíceis
%\end{itemize}
%\end{frame}

%------------------------------------------------

%\begin{frame}
%\frametitle{Duvidas}
%\begin{itemize}
%\item Pós leitura de materiais
%\end{itemize}
%\end{frame}

%------------------------------------------------

\begin{frame}
\Huge{\centerline{OBRIGADO!}}
\end{frame}

%----------------------------------------------------------------------------------------
